\section{Solutions to Exercises from Section 4}
\begin{solution}\label{sol:4.1}\ref{ex:4.1}
    I will also show the cardinality of the powerset since I think it helps with counting all its subsets.
    $$
    \begin{aligned}
        &(1)~\powerset(\{\{ a,b \},\{ c \}\}) = \{ \varnothing,\{ \{ a,b \} \},\{ \{ c \} \},\{ \{ a,b \},\{ c \} \} \}\quad  \lvert \powerset(\{ \{ a,b \},\{ c \} \}) \rvert=2^2=4 \text{ elements}\\\\
        &(2)~\powerset(\{ 1,2,3,4 \})= \{ \varnothing,\{ 1 \},\{ 2 \},\{ 3 \},\{ 4 \},\\
        &\{ 1,2 \},\{ 1,3 \},\{ 1,4 \},\{ 2,3 \},\{ 2,4 \},\\
        &\{ 3,4 \},\{ 1,2,3 \},\{ 1,2,4 \},\{ 1,3,4 \},\{ 2,3,4 \},\{ 1,2,3,4 \} \}=2^4=16 \text{ elements}. \\\\
        &(3)~\powerset(\{ \{ \varnothing \},5 \})= \{ \varnothing,\{ \{ \varnothing \} \},\{ 5 \}, \{ \{ \varnothing \},5 \} \}=2^2=4 \text{ elements}.\\\\
        &(4)~ \powerset(\{ \R,\Q \})= \{ \varnothing, \{ \R \},\{ \Q \},\{ \R,\Q \} \}=2^2=4 \text{ elements}. \\\\
        &(5)~\powerset(\powerset(\{ 2 \}))=\powerset(\{ \varnothing, \{ 2 \} \})=\\
        &\{ \varnothing,\{ \varnothing \}, \{ \{ 2 \} \}, \{ \varnothing, \{ 2 \} \} \}=2^2=4 \text{ elements.}\\\\
        &(6)~\powerset(\{ 1,2 \})\times\powerset(\{ 3 \})=\\
        &\{ \varnothing,\{ 1 \},\{ 2 \},\{ 1,2 \} \}\times \{ \varnothing, \{ 3 \} \}=\\
        & \{ (\varnothing, \varnothing), (\varnothing, \{ 3 \}), (\{ 1 \},\varnothing), (\{ 1 \},\{ 3 \}),\\
        & (\{ 2 \},\varnothing),(\{ 2 \},\{ 3 \}), (\{ 1,2 \},\varnothing),(\{ 1,2 \},\{ 3 \}) \}\\\\
        &(7)~\powerset(\{ a,b \})\times\powerset(\{ 0,1 \})=\\
        & \{ \varnothing,\{ a \},\{ b \},\{ a,b \} \}\times \{ \varnothing,\{ 0 \},\{ 1 \},\{ 0,1 \} \}=\\
        & \{ (\varnothing,\varnothing), (\varnothing, \{ 0 \}),(\varnothing, \{ 1 \}),(\varnothing, \{ 0,1 \}),\\
        &(\{ a \},\varnothing),(\{ a \},\{ 0 \}),(\{ a \},\{ 1 \}),(\{ a \},\{ 0,1 \}),\\
        &(\{ b \},\varnothing),(\{ b \},\{ 0 \}),(\{ b \},\{ 1 \}), (\{ b \},\{ 0,1 \}),\\
        &(\{ a,b \},\varnothing),(\{ a,b \},\{ 0 \}),(\{ a,b \},\{ 1 \}),(\{ a,b \},\{ 0,1 \}) \}\\\\
        &(8)~\powerset(\{ 1,2 \}\times \{ 3 \})= \powerset(\{ (1,3),(2,3) \})=\\
        & \{ \varnothing, \{ (1,3) \},\{ (2,3) \},\{ (1,3),(2,3) \} \}\\\\
        &(9)~\powerset(\{ a,b \}\times \{ 0 \})=\powerset(\{ (a,0),(b,0) \})=\\
        &\{ \varnothing,\{ (a,0) \},\{ (b,0) \},\{ (a,0),(b,0) \} \}
    \end{aligned}
    $$
    We will do the last three examples separately, since there is a bit more to consider. Let's start with $(10)$:
    $$
    (10)~\{ X\in\powerset(\{ 1,2,3 \}): \lvert X \rvert\le 1 \}
    $$
    Notice that we are looking for a \emph{member} $X$, that has size no bigger than $1$. So, let's count the powerset and then see which sets qualify. 
    $$
    \begin{aligned}
        &\powerset(\{ 1,2,3 \})= \{ \varnothing, \{ 1 \},\{ 2 \},\{ 3 \},\\
        &\{ 1,2 \},\{ 1,3 \},\{ 2,3 \},\{ 1,2,3 \} \}=2^3=8 \text{elements}
    \end{aligned}
    $$
    Now, for the condition $X\in\powerset(\{ 1,2,3 \})$. We have 
    $$
    \{\varnothing, \{ 1 \},\{ 2 \},\{ 3 \}\}.
    $$
These are all the singletons and the empty set. For $(11)$, we have 
$$
\{ \varnothing,\{ \varnothing \}, \{ \{ 1 \} \},\{ \{ 2 \} \},\{ \{ 3 \} \}, \{ \{ 1,2 \} \},\{ \{ 1,3 \} \},\{ \{ 2,3 \} \},\{ \{ 1,2,3 \} \} \}.
$$
This might seem strange, but these are all of cardinality at most than $1$. Notice that the subsets of the powerset are all enclosed in an extra pair of curly braces, so they are counted as having one element. Lastly, for $(12)$ the condition changes slightly, we need to find $X\in\powerset(\{ 1,2,3 \})$ such that $2\in X$. So, 
$$
\{\{ 2 \},\{ 1,2 \},\{ 2,3 \},\{ 1,2,3 \}\}.
$$    
\end{solution}

\begin{solution}\label{sol:4.2}\ref{ex:4.2}
    We let $\lvert A \rvert= m ,\lvert B \rvert= n$, then 
    $$
    \begin{aligned}
        &(1)~\lvert \powerset\powerset\powerset(A) \rvert=2^{2^{2^{m}}}\\\\
        &(2)~\lvert \powerset\powerset(A) \rvert=2^{2^{m}}\\\\
        &(3)~\lvert \powerset(A\times B) \rvert=2^{m\cdot n}\\\\
        &(4)~\lvert \powerset(A)\times \powerset(B) \rvert=2^{m}\cdot 2^{n}\\\\
        &(5)~\lvert \{ X\in\powerset(A):\lvert X \rvert\le 1 \} \rvert= m+1\\\\
        &(6)~\lvert \powerset(A\times\powerset(B)) \rvert=2^{m\cdot 2^{n}}\\\\
        &(7)~ \lvert \powerset\powerset\powerset(A\times\varnothing) \rvert=4\\\\
        &(8)~\lvert \{ X\subseteq\powerset(A):\lvert X \rvert\le 1 \} \rvert=2^{m}+1.
    \end{aligned}
    $$
    \emph{Some Further Comments:} For $(5)$ it is important to understand the following. We are looking at members of the powerset, as in the previous exercise, we saw that these would be all the singletons and the empty set. Since we have exactly $m$ singletons and we add the empty set to that, so we get $m+1$. 
    Since the product of any set with the empty set is empty, we have to calculate $\lvert \powerset\powerset\powerset(\varnothing) \rvert$. The idea in $(8)$ is the same as in the previous exercise. 
\end{solution}